\documentclass[10pt,letterpaper]{article}
\usepackage{cornell-notes}

\title{Chapter 21: Power-Series Composition}
\author{}
\date{}

\setDocTitle{Chapter 21: Power-Series Composition}
\setDocSubtitle{Combinatorial Algorithms}
\setDocVersion{v1.0}
\setDocStatus{Study Companion}
\setDocOwner{Combinatorial Algorithms Cornell Notes}
\setDocAuthor{}
\setDocDate{\today}
\setDocSummary{Cornell-format study and review notes for Chapter 21: Power-Series Composition.}

\setCornellCollection{Combinatorial Algorithms Cornell Notes}
\setCornellUnitType{Chapter}
\setCornellUnitNumber{21}
\setCornellUnitTitle{Power-Series Composition}
\setCornellCompanionLabel{Study and review companion}

\hypersetup{pdftitle={Chapter 21: Power-Series Composition},pdfsubject={Cornell notes},pdfauthor={}}

\begin{document}
\maketitle
\begin{tcolorbox}[colback=Paper,colframe=Ink]{\small\bfseries\color{Teal}CORNELL NOTES}\par{\LARGE\bfseries Chapter 21: Composition of Power Series (POWSER)}\par\medskip\textbf{Topic:} Truncated formal-series transforms\hfill\textbf{Date:} \today\par\textbf{Study purpose:} Compute coefficients by specialized recurrences and inverse normalization.\end{tcolorbox}
\begin{tcolorbox}[colback=Teal!5,colframe=Teal,title=Essential question]Which coefficient recurrences make important power-series compositions computable in quadratic rather than combinatorial work?\end{tcolorbox}
\begin{tcolorbox}[colback=white,colframe=Mist,title=Learning objectives]\begin{itemize}\item State the zero-constant-term conditions for composition.\item Derive the recurrence for $(1+h)^\alpha-1$.\item Specialize to $\exp(h)-1$.\item Explain general composition through the first nonzero term of $h$.\item Describe the compositional-inverse method.\item Select among POWSER's four options.\end{itemize}\textbf{Prerequisites:} formal power series, differentiation, coefficient comparison, composition.\end{tcolorbox}
\section*{Cornell notes}
\begin{longtable}{@{}Q!{\color{Mist}\vrule width 1pt}A@{}}\cornellhead
\cue{What is represented?}{Arrays store the first $n$ nonconstant coefficients of series such as $f(z)$, $g(z)$, and $h(z)$. Composition requires $h(0)=0$ so each requested output coefficient depends on finitely many input coefficients.}
\cue{Option 1: what identity is used?}{For $F(z)=(1+h(z))^\alpha-1$, differentiate and multiply by $1+h$:
\[
F'(1+h)=\alpha(1+F)h'.
\]
Equating coefficients yields a triangular recurrence: coefficient $a_j$ depends only on $c_1,\ldots,c_j$ and earlier $a_i$. All $n$ coefficients require $O(n^2)$ arithmetic.}
\cue{Option 2: why is exponential simpler?}{For $F=\exp(h)-1$, logarithmic differentiation gives $F'=(1+F)h'$. This is the same recurrence form with specialized parameters and computes the first $n$ coefficients in $O(n^2)$ work.}
\cue{Option 3: general $g(h)$}{Let the first nonzero term of $h$ be $c_qz^q$. Factor $h(z)=c_qz^q(1+u(z))$. For each nonzero outer coefficient $b_m$, compute only the needed power $(1+u)^m$ with Option 1, shift degrees by $mq$, scale by $b_mc_q^m$, and accumulate. Powers whose outer coefficients are zero are skipped.}
\cue{Why not multiply powers successively?}{Building $h,h^2,h^3,\ldots$ uses two work arrays and computes unneeded lower powers when only a high power is required. It also propagates rounding error from each power to the next. The selected recurrence computes each required power independently.}
\cue{Option 4: inverse composition}{To compute $g=f\circ h^{-1}$ when $h'(0)\ne0$, successively compose $h$ with simple inverse factors until it agrees with $z$ through degree $n$. Apply the same coefficient transformations to $f$. The transformed $f$ coefficients are then those of $g$.}
\cue{Unified recurrence}{Options 1–3 share one parameterized triangular formula controlling exponent, degree shift, normalization, and scale. Option 4 uses nested in-place Taylor-like transforms on both the inner series and the outer coefficients.}
\cue{Cost and storage}{The option table gives $O(n^2)$ operation time for all four supported tasks. Arrays $A,B,C$ hold result and inputs; $D$ is length-$n$ workspace. Option 4 requires a nonzero linear coefficient $C(1)$.}
\cue{Cautions}{All series are truncated, so state the requested degree. Division by powers of the leading coefficient can magnify error. Composition and inverse composition fail when the required zero constant term or nonzero derivative condition is violated.}
\cue{Retrieval practice}{\begin{enumerate}\item Why must $h(0)=0$?\item What identity drives Option 1?\item What changes for the exponential?\item What does $q$ denote in Option 3?\item Why must $h'(0)\ne0$ for Option 4?\end{enumerate}}
\end{longtable}
\begin{tcolorbox}[colback=Ink!4,colframe=Ink,title=Cornell summary]
POWSER replaces direct combinatorial coefficient formulas with triangular recurrences. Differentiating $(1+h)^\alpha-1$ and $\exp(h)-1$ produces recurrences in which coefficient $j$ depends only on earlier results, giving quadratic work. A general composition $g(h)$ factors out the first nonzero term of $h$ and independently computes only the powers requested by nonzero outer coefficients. Inverse composition normalizes $h$ to the identity through successive simple inverse factors while applying the same transforms to $f$. The four options share array conventions and $O(n^2)$ truncation cost, but require careful attention to zero constant terms, the leading coefficient, and numeric stability.
\end{tcolorbox}
\end{document}
